\noindent\textit{2008 manuscript.}

\section{Kinetic Theory}

\noindent\textit{The wall-collision argument for \(P=N\tau/V\) is in \S\ref{sec:dilute-gas}.}

\subsection*{Kinetic Theory of the Ideal Gas Law}

Consider molecule strike unit area of wall. \\
Let $v_z \equiv $ velocity component normal to plane of wall. \\
Suppose molecules, of mass $M$, reflected specularly (mirror-like) from wall,
\[
\Delta p_z = -2M|v_z|
\]
Let $a(v_z)dv_z$, number of molecules per unit volume with $z$-component of velocity between $v_z$ and $v_z + dv_z$.  
\[
\int a(v_z)dv_z = \frac{N}{V} = n
\]
$a(v_z) v_z dv_z$ number of molecules in $(v_z, v_z + dv_z)$  velocity range that strike unit area of wall in (per) unit time
\[
\text{pressure } p = \int_0^{\infty}2M v_za(v_z)v_z dv_z = M\int_{-\infty}^{\infty}v_z^2 a(v_z) dv_z = Mn\langle v_z^2 \rangle
\]
$\frac{1}{2}M\langle v_z^2 \rangle = \frac{1}{2} \tau$ by equipartition of energy (Ch.3)
\[
p = nM \langle v_z^2 \rangle = n\tau = \frac{N\tau}{V}
\]
\subsubsection*{Maxwell Distribution of Velocities}  

cf. Ch.6. distribution function of ideal gas $f(\epsilon_n) = \lambda \exp{ \left( \frac{-\epsilon_n}{\tau} \right)}$

Recall, Ch. 6, Sec. ``Classical Limit'' of Kittel and Kroemer \cite{CKittelHKroemer1980}, \textbf{an ideal gas is defined as a system of free noninteracting particles in the classical regime}.  \\
$f(\epsilon) \equiv $ average occupancy of an orbital at energy $\epsilon$ \\
$\epsilon \equiv $ energy of orbital occupied by 1 particle; not energy of system of $N$ particles \\
Fermi-Dirac and Bose-Einstein distribution $f(\epsilon) = \frac{1}{ \exp{ [ (\epsilon - \mu )/\tau ] } \pm 1 }$ \\
In order for $f(\epsilon) \ll 1$ $\forall \, $ orbitals, $\exp{ [ (\epsilon-\mu)/\tau ] } \gg 1\, \forall \, \epsilon$.  
\[
\Longrightarrow f(\epsilon) \simeq \exp{ [(\mu - \epsilon )/\tau]} = \lambda \exp{ (-\epsilon/\tau)} \quad \, \lambda \equiv \exp{ \left( \frac{ \mu }{\tau} \right) }
\]
$f(\epsilon)$, average occupancy of orbital of energy $\epsilon$, is classical distribution function.  

particle in a box: $\epsilon_n = \frac{1}{2M} \left( \frac{\pi n}{L} \right)^2 $ (for, recall $\frac{1}{2M} \left( \frac{1}{i} \partial \right)^2 \psi = \frac{-1}{2M} \partial^2 \psi = E \psi$)

number of orbitals in range of quantum number $(n,n+dn)$, probability such orbital is occupied

\[
(\frac{1}{2} \pi n^2 dn)f(\epsilon_n) = \frac{1}{2} \pi n^2 \lambda \exp{ (-\epsilon_n/\tau) } dn
\]
\[
\frac{1}{2} Mv^2 = \frac{1}{2M} \left( \frac{ \pi n}{ L } \right)^2 \text{ or } n^2 = \frac{ (ML)^2 }{ \pi^2} v^2 \text{ or } n = \frac{MLv}{\pi}
\]

Consider system of $N$ particles in volume $V$.  

Let $NP(v) dv$ number of atoms with velocity magnitude in range $dv$ at $v$
\[
NP(v) dv = \frac{1}{2} \pi n^2 \lambda \exp{ ( -\epsilon_n /\tau)} \frac{dn}{dv} dv = \frac{1}{2} \pi \lambda \left( \frac{ ML }{ \pi } \right)^3 v^2 \exp{ \left( \frac{ -Mv^2 }{2\tau } \right) }dv
\]
cf. Ch. 6, $\lambda = \frac{n}{n_Q} = \frac{N}{L^3} \left( \frac{ 2\pi \hbar^2 }{ M \tau } \right)^{3/2}$

\[
\frac{1}{2} \pi \frac{N}{L^3} \left( \frac{2\pi }{ M \tau} \right)^{3/2} \left( \frac{ ML }{\pi } \right)^3 = 4\pi N \left( \frac{N}{2\pi \tau} \right)^{3/2}
\]
\begin{correction}{Maxwell prefactor}
The right-hand side above is not the prefactor of the speed density in the next display.
That density is normalized, as checked in \S\ref{sec:maxwell-norm}.
\end{correction}
\begin{equation}
\Longrightarrow P(v) = 4\pi \left( \frac{M}{2\pi \tau } \right)^{3/2} v^2 \exp{ \left( \frac{- Mv^2}{ 2\tau } \right) } 
\end{equation}
$P(v)$ is \textbf{Maxwell velocity distribution}, $P(v)dv$ is probability particle has speed in $dv$ at $v$

\subsubsection*{Experimental verification}

velocity distribution of atoms which exit from slit of oven.  

exit beam weighted in favor of atoms of high velocity at expense of those at low velocity. 

weight factor is velocity component $v\cos{\theta}$ normal to plane of hole. 

\[
\begin{gathered}
  \int (\cos{\theta}) dr rd\varphi r \sin{\theta} d\theta = (\frac{1}{3} (2\pi ) R^3) \int \cos{\theta} \sin{\theta} d\theta = ( \frac{2\pi}{3} R^3) \int \frac{\sin{2\theta} d\theta}{ 2 } = \frac{2\pi R^3}{3} \left. \left( \frac{ - \cos{2\theta} }{4} \right) \right|_0^{\pi/2} = \\
  = \frac{ 2\pi R^3}{3} \left( \frac{1+ 1}{4} \right) = \frac{ 2\pi R^3}{3} \left( \frac{1}{2} \right)
\end{gathered}
\]
Probability atom leaves hole will have velocity in $(v,v+dv)$ : $P_{\text{beam}}(v) dv$

\[
P_{\text{beam}}(v) \propto vP_{\text{Maxwell} } \propto v^3 \exp{ \left( \frac{ -Mv^2}{ 2\tau } \right) }
\]
with, recall $P_{\text{Maxwell}} = 4\pi \left( \frac{ M }{ 2\pi \tau } \right)^{3/2} v^2 \exp{ \left( \frac{-Mv^2}{ 2\tau } \right) }$ 

$P_{\text{beam}}$ distribution of transmission through a hole is Maxwell transmission distribution. 

\[
\begin{gathered}
  \langle v_{\text{out}} \rangle = \int \int_0^{\pi/2} v\cos{\theta} \sin{\theta} d\theta P_{\text{Maxwell}}(v) dv = 4\pi \left( \frac{ M}{ 2\pi \tau} \right)^{3/2} \frac{1}{2} \int_0^{\infty} v^3 \exp{ \left( \frac{-M}{2\tau} v^2 \right) } dv = \\
  = 4\pi \left( \frac{M}{2\pi \tau} \right)^{3/2} \frac{1}{2} \left( \frac{1}{-2\alpha} \right) \left[ 0 - \frac{1}{\alpha} \right] = 4\pi \left( \frac{ M}{2\pi \tau} \right)^{3/2} \frac{1}{4} \left( \frac{4\tau^2}{M^2} \right) = \boxed{ \frac{ 2^{1/2} }{ \pi^{1/2} } \left( \frac{\tau}{M} \right)^{1/2} }
\end{gathered}
\]
for (doing the integration by hand) 
\[
\begin{aligned}
  & (e^{-\alpha v^2} )' = -2\alpha v e^{-\alpha v^2} \\
  & (v^2 e^{-\alpha v^2} )' = -2\alpha v^3 e^{-\alpha v^2} + 2v e^{-\alpha v^2} \\ 
  & (v^2 e^{-\alpha v^2} + \frac{e^{-\alpha v^2} }{ \alpha } )' = -2\alpha v^3 e^{-\alpha v^2}
\end{aligned}
\]




\subsubsection*{Particle Diffusion}

pp. 399 of Kittel and Kroemer \cite{CKittelHKroemer1980}

Consider a system.  \\
One end in diffusive contact with reservoir at chemical potential $\mu_1$ \\
Other end in diffusive contact with reservoir at chemical potential $\mu_2$ \\
Constant temperature $\tau$.  \\
If $\mu_1 > \mu_2$, particle flow through system from reservoir 1 to reservoir 2; $1\to 2$.   \\
$n_i \equiv $ particle concentration in $i$ 

Take
\begin{equation}
\mathbf{j}_n = -D \text{grad}n
\end{equation}
which is \textbf{Fick's law}, and where $D \equiv$ particle diffusion constant or \textbf{diffusivity}.  

Mean free path $l$.  Particles freely travel over $l$.  

Assume in a collision at $z$, particles come into local equilibrium at local chemical potential $\mu(z)$, local concentration $n(z)$.  

At $z$, particle flux density in positive $z$ direction \qquad $\frac{1}{2} n(z-l_z) \overline{c}_z$ \\
\phantom{At $z$, } particle flux density in negative $z$ direction \qquad $-\frac{1}{2} n(z+l_z) \overline{c}_z$

Note $n(z-l_z)$ is particle concentration at $z-l_z$
\[
J_n^z = \frac{1}{2} \left[ n(z-l_z) - n(z+l_z) \right] \overline{c}_z = -\frac{dn}{dz} \overline{c}_z l_z
\]
where $\begin{aligned} & \quad \\
  & \overline{c}_z = \overline{c} \cos{\theta}  \\
  & \overline{l}_z = \overline{l} \cos{\theta}  
\end{aligned}$

\[
\langle \overline{c}_z l_z \rangle = \overline{c} \overline{l} \frac{ \int_{\text{hemisphere} } \cos^2{\theta} dS }{ \int_{\text{hemisphere} } dS } = \overline{c}\overline{l} \frac{ \int_0^{\pi/2}d\theta \int_0^{2\pi } d\varphi \cos^2{\theta} \sin{\theta} d\theta }{ \int_0^{\pi/2} d\theta \sin{\theta} \int_0^{2\pi } d\varphi } = \overline{c} \overline{l} \frac{1}{3}
\]
Comparing with Fick's law, 
\[
J_n^z = \frac{-1}{3} \overline{c}\overline{l} \frac{dn}{dz} \text{ or } \mathbf{J}_N = -\frac{1}{3} \overline{c}\overline{l} \nabla n
\]
For diffusivity $D$ is then $D = \frac{1}{3} \overline{c}\overline{l}$.  

Now recall that $\overline{l}$, the \emph{mean free path}, was derived from kinetic theory:
\[
l = \frac{1}{n\pi d^2}
\]
where $d$ is the diameter of the particle.  

The mean thermal velocity was derived from the Maxwell distribution:
\[
\overline{c} = \left( \frac{ 8 \tau }{M \pi} \right)^{1/2}
\]





\subsection*{Problems}

\problemhead{1. Mean speeds in a Maxwellian distribution}

\begin{enumerate}
  \item[(a)] $P(v) = 4\pi \left( \frac{M}{2\pi \tau} \right)^{3/2} v^2 \exp{ \left( \frac{-Mv^2}{2\tau} \right) }$

Now
\[
\begin{gathered}
  \int_0^{\infty} dv v^4 \exp{ (-\alpha v^2 )} = \int_0^{\infty} v^3 \left( \frac{ \exp{ (-\alpha v^2) } }{ -2 \alpha } \right)' = 0 - \int 3v^2 \frac{ \exp{ (-\alpha v^2)} }{ -2\alpha } = \int_0^{\infty} \frac{3}{2\alpha} v^2 \exp{ (-\alpha v^2 )} = \\
  = \frac{3}{2\alpha } \int v \left( \frac{e^{-\alpha v^2 }}{ -2\alpha } \right)' = \frac{3}{2\alpha } \left[ 0 - \int \frac{e^{-\alpha v^2 } }{ -2\alpha } \right] = \frac{3}{4\alpha^2 } \frac{ \sqrt{ \pi }}{ 2\sqrt{ \alpha }}
\end{gathered}
\]
So
\[
\langle v^2 \rangle = \int_0^{\infty} v^2 P(v) = 4\pi \left( \frac{M}{2\pi \tau } \right)^{3/2} \frac{3}{4 \left( \frac{M}{2\tau} \right)^2 } \frac{ \sqrt{\pi }}{ \sqrt{ \frac{M}{2\tau} } } \frac{1}{2} = \frac{3}{ \frac{M}{2\tau} } \frac{1}{2} = \frac{3\tau }{M}
\]
so 
\[
\boxed{ v_{rms} = \left( \frac{3\tau}{M} \right)^{1/2} }
\]
Now
\[
\langle v^2 \rangle = 3\langle v_x^2 \rangle = \frac{3\tau}{M} \text{ so } \langle v_x^2 \rangle^{1/2} = \left( \frac{\tau}{M} \right)^{1/2}
\]
  \item[(b)] \[
\begin{gathered}
  \frac{ \partial P(v) }{ \partial v} = \frac{2 P(v)}{v} + P(v) \left( \frac{-M}{\tau} \right)v = P(v) \left[ \frac{2}{v} - \frac{M}{\tau}v \right] = 0  \text{ if } v_{mp} = 0 \text{ or } v_{mp} = \boxed{ \sqrt{ \frac{2\tau}{M} } }
\end{gathered}
\]
where $mp$ stands for most probable.  

\[
v_{mp} = \sqrt{ \frac{2\tau}{M} } < \sqrt{ \frac{3\tau}{M} } = v_{rms}
\]
  \item[(c)] \[
\begin{gathered}
  \overline{c} = \int_0^{\infty} dv v P(v) =: \langle v \rangle = \int_0^{\infty}dv 4\pi \left( \frac{M}{ 2\pi \tau} \right)^{3/2} v^3 \exp{ \left( \frac{-Mv^2}{2\tau} \right) } = 4 \pi \left( \frac{M}{2\pi \tau} \right)^{3/2} \int_0^{\infty}dv v^2 \left( \frac{ e^{ -\frac{Mv^2}{2\tau} } }{ \frac{-M}{\tau} } \right)' = \\
  = 4\pi \left( \frac{M}{2\pi \tau} \right)^{3/2} \left[ 0 - \int_0^{\infty} dv 2v \frac{e^{ -\frac{Mv^2}{2\tau } } }{ \frac{-M}{\tau} } \right] = 4\pi \left( \frac{M}{2\pi \tau} \right)^{3/2} \left[ \frac{2\tau}{M} \left. \left( \frac{e^{-\frac{Mv^2}{ 2\tau } } }{ -\frac{M}{\tau} } \right) \right|_0^{\infty} \right] = 4 \pi \left( \frac{M}{2\pi \tau} \right)^{3/2} \frac{2\tau^2}{M^2} = \\
   = \frac{1}{(2\pi )^{3/2} } 8 \pi \sqrt{ \frac{\tau}{M}} = \boxed{ \sqrt{ \frac{8\tau}{\pi M} } }
\end{gathered}
\]
Note
\[
\begin{gathered}
  \frac{v_{rms}}{ \overline{c}} = \frac{ \left( \frac{3\tau}{M} \right)^{1/2} }{ \left( \frac{8\tau}{\pi M} \right)^{1/2} } = \left( \frac{3\pi }{8} \right)^{1/2}
\end{gathered}
\]
\item[(d)] $|v_z| = |v\cos{\theta} | = v|\cos{\theta}|$.  For for a fixed $\theta$, $v\cos{\theta}$ is 1 point on a sphere, whereas $v|\cos{\theta}|$ could be 2 pts on a sphere.  

Now
\[
\begin{gathered}
\int_0^{\pi} |\cos{\theta}| \sin{\theta}d\theta = \int_0^{\pi/2} \cos{\theta} \sin{\theta} d\theta + \int_{\pi/2}^{\pi} - \cos{\theta} \sin{\theta} d\theta = \int_0^{\pi/2} \frac{\sin{2\theta} d\theta}{ 2} + \int_{\pi/2}^{\pi} \frac{ - \sin{2\theta} d\theta}{2} = \\
= \left. \left( \frac{ \cos{2\theta} d\theta}{-4} \right) \right|_0^{\pi/2} + \left. \left( \frac{ \cos{2\theta} }{ 4} \right) \right|_{\pi/2}^{\pi} = \frac{1}{2} + \frac{1 + 1 }{4} = 1 
\end{gathered}
\]
$P(v)dv$ is probability that particle has speed or velocity magnitude in $(v,v+dv)$.  \\
It counts the event of $v\cos{\theta}$ and $v\cos{(\pi-\theta)}$, so it counts $v|\cos{\theta}|$ twice.  $P(v)dv = \int_0^{\pi/2} d\theta 2 P(v|\cos{\theta}|) $
\[
\Longrightarrow \overline{c}_z = \langle |v_z| \rangle = \langle |v\cos{\theta}| \rangle = \langle v |\cos{\theta} | \rangle = \int_0^{\pi} \int_0^{\infty} v \frac{ P(v) dv}{2} |\cos{\theta}| d\theta = \left( \frac{2\tau}{\pi M } \right)^{1/2}
\]

\end{enumerate}



\section{Propagation}

\subsection*{Heat Conduction Equation}

nonrelativistic case: 

Let manifold $N = \mathbb{R} \times M$, with $\text{dim}M = n$.  

Let $\mathbf{J} \in \mathfrak{X}(N) = \mathfrak{X}(\mathbb{R}\times M)$ be a vector field in $N$.  \\
Let $\begin{aligned} & \quad \\ 
  & \rho \in C^{\infty}(N) \\ 
  & \rho = \rho(t,x) \text{ locally } \end{aligned}$ be a smooth function on $N$.  \\
Let $J \in \Omega^1(N)$ be a 1-form on $N$ that is isomorphic to $\mathbf{J}$ (Tangent-Cotangent isomorphism theorem), i.e.
\[
\begin{aligned}
& J = \mathbf{J}^{\flat}  \\
& J_i = g_{ij} \mathbf{J}^j
\end{aligned}
\]
with $g_{ij}$ being the metric on $N$ (not just $M$!)

Note that as $N = \mathbb{R}\times M$, $g_{0j} = \delta_{0j}$

The local form of $\mathbf{J}$ is the following:
\[
\mathbf{J} = \rho \frac{ \partial }{ \partial t} + \mathbf{j}^i \frac{ \partial }{ \partial x^i} \quad \quad \, i=1\dots n
\]

So 
\[
\begin{gathered}
  J = J_i dx^i = g_{ij} \mathbf{J}^j dx^i = \rho dt + j_k dx^k \quad \quad \, k = 1 \dots n 
\end{gathered}
\]
Thus
\begin{equation}
  J = \rho dt + j_k dx^k \quad \quad \, k=1\dots n
\end{equation}

Now do the Hodge star operator, resulting in a $n$-form
\[
*J \in \Omega^n(N)
\]
and so
\[
\begin{gathered}
  *J = \rho * dt + j_k * dx^k = \rho \text{vol}^n + i_{\mathbf{j}}\text{vol}^{n+1}
\end{gathered}
\]
as
\[
\begin{gathered}
  \rho * dt = \rho \frac{ \sqrt{g}}{ n !} \epsilon^0_{ \, \, i_1 \dots i_n} dx^i_1 \wedge \dots \wedge dx^{i_n} \quad \quad \, i_1 \dots i_n \in \lbrace 1 \dots n \rbrace \\ 
  j_k * dx^k = g_{kl} j^l \frac{ \sqrt{g}}{ (n+1)!} \epsilon^k_{ \, \, i_1 \dots i_n} dx^{i_1} \wedge \dots \wedge dx^{i_n} = i_{\mathbf{j}} \text{vol}^{n+1} 
\end{gathered}
\]
so thus
\[
*J = \rho \text{vol}^n + i_{\mathbf{j}} \text{vol}^{n+1}
\]

Hence
\begin{equation}
  d*J = \frac{ \partial }{ \partial t} (\rho \sqrt{g} ) \frac{1}{ \sqrt{g}} \text{vol}^{n+1} + \frac{ \partial }{ \partial x^k} (\sqrt{g} j_k) \frac{1}{ \sqrt{g}} \text{vol}^{n+1} = d(\rho \text{vol}^n) + di_{\mathbf{j}} \text{vol}^{n+1}
\end{equation}

Special case: $\frac{ \partial \sqrt{g}}{ \partial t} = 0$

\[
d*J = \frac{ \partial \rho}{ \partial t} \text{vol}^{n+1} + \left( \frac{ \partial }{ \partial x^k} \ln{ \sqrt{g}} \right) j_k \text{vol}^{n+1} + \frac{ \partial j_k}{ \partial x^k} \text{vol}^{n+1} = 0 
\]
\[
\Longrightarrow \frac{ \partial \rho }{ \partial t} + \left( \frac{ \partial }{ \partial x^k} \ln{ \sqrt{g}} \right) j_k + \frac{ \partial j_k }{ \partial x^k } = 0 
\]

Special case: if $\sqrt{g}$ constant, $\frac{ \partial \rho }{ \partial t } + \frac{ \partial j_k}{ \partial x^k} = 0 $

Let $j \equiv -D \mathbf{d}\rho$ ($j$ is a closed form on $M$) where $\mathbf{d}\rho = \frac{ \partial \rho }{ \partial x^i } dx^i$ \quad \, $i=1 \dots n$, $D$ constant

\[
d * J  = d\rho \text{vol}^n + - Dd * \text{d}\rho = 0 
\]
For flat metric, $\sqrt{g}$ constant,
\[
\frac{ \partial \rho }{ \partial t} = D \frac{ \partial^2 \rho }{ ( \partial x^k )^2}
\]

For relativistic case, 

Consider manifold $M$, with dimension $\text{dim}M = n+1$
\[
\mathbf{J} = \rho \frac{ \partial }{ \partial t} + \mathbf{j}^i \frac{ \partial }{ \partial x^i}
\]

Let $J = \mathbf{J}^{\flat}$.  $J_{\mu} = g_{\mu \nu} \mathbf{J}^{\nu}$

For special case of flat Minkowski space,
\[
J  = - \rho dt + j_i dx^i
\]

\[
\begin{aligned}
  *J & = - \rho * dt + j_i * dx^i \\ 
  & -\rho \frac{ \sqrt{g}}{ (n+1)!} \epsilon^0_{ \, \, i_1 \dots i_n} dx^{i_1} \wedge \dots \wedge dx^{i_n} + j_i \frac{ \sqrt{g} }{ (n+1)!} \epsilon^i_{ \, \, \mu_1 \dots \mu_n } dx^{\mu_1} \wedge \dots \wedge dx^{\mu_n}
\end{aligned}
\]
\[
d* J = -\frac{ \partial (\rho \sqrt{g}  )}{  \partial t } \frac{1}{ \sqrt{g} } \text{vol}^{n+1} + \frac{ \partial (j_i \sqrt{g} )}{  \partial x^i } \frac{1}{ \sqrt{g}} \text{vol}^{n+1} = 0 
\]
\[
\Longrightarrow - \frac{ \partial (\rho \sqrt{g})}{  \partial t} + \frac{ \partial (j_i \sqrt{g} )}{ \partial x^i } = 0 
\]

Fick law (14.19) for particle flux density, $j = -D_n \mathbf{d}n$ where $\begin{aligned}  & \quad \\ 
  & D_n  & \text{ particle diffusivity constant } \\ 
  & n   & \text{ particle concentration } \end{aligned}$ 

\[
J = n dt + j
\]

thermal conductivity; homogeneous medium
$\widehat{C}$ heat capacity per unit volume.  $j_u = -K \mathbf{d}\tau$

$J = \widehat{C} \tau dt + j_u$

\begin{equation}
\widehat{C} \frac{ \partial \tau}{ \partial t} + \frac{ \partial (j_u)_k}{ \partial x^k} = 0 \quad \quad \quad \, (5)
\end{equation}

\begin{equation}
  \frac{ \partial \tau}{ \partial t} = D_{\tau} \frac{ \partial^2 \tau}{ ( \partial x^k)^2 } \quad \quad \, D_{\tau} \equiv K/\widehat{C} \quad \quad \quad \, (6) 
\end{equation}

\subsection*{Propagation of Sound Waves in Gases}

pressure associated with sound wave 

\begin{equation}
  \delta p = \delta p_0 \exp{ [ i (kx - \omega t) ] } \quad \quad \quad \, (27) 
\end{equation}

Suppose ideal gas: 
\begin{equation}
  pV = N \tau \text{ or } p = \rho \tau /M \quad \quad \quad \, (28)
\end{equation}
Consider ``solid ball'' or ``billiard ball'' particle (extended particle, not pt. particle, but no internal structure)
\[
\rho = \frac{NM}{V}
\]
Force on particle 
\[
F = \frac{dP}{dt} = \frac{d}{dt} \int \rho \text{vol}^n u = \frac{M}{V} \int \mathcal{L}_{ \frac{ \partial }{ \partial t} + \mathbf{u} }\mathbf{u} = \frac{M}{V} \int \frac{ \partial u}{ \partial t} + [u,u]
\]
Suppose $[u,u]=0$ (certainly for flat spaces; what about for curved spaces? $[u,u] \neq 0$? Possibly? I don't know. EY: 20150317

\begin{equation}
dU + p dV = \tau d\sigma
\end{equation}

define fractional deviations $s,\theta$

\begin{equation}
  \begin{aligned}
    & \rho = \rho_0(1+s) \\
    & \tau = \tau_0 (1 + \theta) 
\end{aligned} \quad \quad \quad \, (5)
\end{equation}
where $\rho_0$, $\tau_0$ are density and temperature in absence of sound wave. \\
assume $u,s,\theta$ have form of traveling $\exp{ [ i (kx - \omega t)]}$

\begin{align}
& \omega u - \left( \frac{k \tau}{M} \right) s - \left( \frac{ k \tau }{M} \right) \theta =0  \quad \quad \quad \, (39) \\ 
& \omega s - ku =0 \quad \quad \quad \, (40) \\ 
& \tau \widehat{C}_V \theta - ps = 0 \text{ or } \widehat{C}_V \theta - ns = 0 \quad \quad \quad \, (41)
\end{align}

\[
\begin{gathered}
  \omega u = \frac{k\tau}{ M} \left( 1 + \frac{n}{ \widehat{C}_V} \right) s \\ 
  \omega = \frac{k\tau}{M} \left( 1 + \frac{n}{ \widehat{C}_V} \right) \frac{k}{\omega}
\end{gathered}
\]
So 
\begin{equation}
  \omega = \left( \frac{ \gamma \tau}{M} \right)^{1/2} k  \quad \quad \quad \, (42)
\end{equation}

\[
\begin{gathered}
  \gamma = \frac{ \widehat{C}_V + n }{ \widehat{C}_V } = \frac{\widehat{C}_p}{ \widehat{C}_V } \\ 
 v_s = \frac{ \partial \omega}{ \partial k} = \left( \frac{\gamma \tau}{ M} \right)^{1/2}
\end{gathered}
\]

\subsection*{Problems}

\problemhead{1} \textbf{Fourier analysis of pulse}

$t=0$ 

\begin{equation}
  \theta(x,0) = \delta(x) = \frac{1}{2\pi} \int_{-\infty}^{\infty} dk \exp{ (ikx)} \quad \quad \quad \, (58) 
\end{equation}

\begin{equation}
  \theta(x,t) = \frac{1}{2\pi} \int_{-\infty}^{\infty} dk \exp{ [ i (kx-\omega t) ]} \quad \quad \quad \, (59) 
\end{equation}

Given a dispersion relation at this form:
\begin{equation}
  Dk^2  = i \omega \quad \quad \, (10)
\end{equation}

\begin{equation}
\theta(x,t) = \frac{1}{2\pi} \int_{-\infty}^{\infty} dk \exp{ [ ikx - Dk^2 t] } \quad \quad \quad \, (60)
\end{equation}

and so, doing the Gaussian integral,
\begin{equation}
  \theta(x,t) = \frac{1}{ \sqrt{ 4\pi Dt} } \exp{ \left( \frac{-x^2}{ 4Dt } \right)} \quad \quad \quad \, (14)
\end{equation}


\problemhead{Diffusion in two and three dimensions}

\begin{enumerate}
\item[(a)] \[
\begin{aligned}
  & \frac{ \partial \theta_2}{ \partial t} = - \frac{ \theta_2}{t} + \frac{r^2}{4} \frac{ \theta_2}{ Dt^2 } \\ 
  & \frac{ \partial^2\theta_2}{ (\partial x^i)^2} = - \frac{1}{2} \frac{ \theta_2}{Dt} - \frac{1}{2} \frac{\theta_2}{ Dt} + \frac{1}{4} \frac{ \theta_2}{ D^2 t^2} (x^2+ y^2) 
\end{aligned}
\]
\[
\Longrightarrow \frac{ \partial \theta_2}{ \partial t} = D \frac{ \partial^2 \theta_2}{ (\partial x^i)^2 }
\]
\item[(b)]
\item[(c)]
\end{enumerate}

